\section{Limitations and conclusion}
\label{sec:conclusion}

\runin{Limitations.} The corpus is a machine transcript and a finding aid. It recovered 76.1\% of test reference words, and 49.3\% in the EOB, so every quotation should be checked against the audio, to which every word is linked. FRUS editors chose the test excerpts; a test set drawn at random from the corpus would measure the whole record. Correction depends on a proprietary model, but its edits are marked and reversible. Voices without a class of their own were sometimes given a trained speaker's name (16.9\% of their test words). Withdrawn passages stay muted, and passages about private persons call for care in republication.

\runin{Conclusion.} Two recognizers that err differently, reconciled against the archive's own records, turn the released Nixon tapes into a searchable, speaker-labeled text of 19.0 million words. On held-out conversations this text recovered 76.1\% of words and 80.1\% of names, against 69.2\% and 60.0\% for Whisper alone. It tells a historian where to listen, with measured accuracy, word timings and speaker labels that neither the subject logs nor earlier machine transcripts provide.
